C++17 中另一个实用的包装类型是 \texttt{std::variant}。它是一种类型安全的联合体（union）——你可以存储不同的类型变体，并获得正确的对象生命期保证。这种新类型相比 C 风格的 union 有着巨大的优势。你可以把各种类型都放进去——无论是像 \texttt{int}、\texttt{float} 这样简单的类型，还是像 \texttt{std::vector<std::string>} 这样复杂的实体。在所有情况下，对象都会被正确地初始化和清理。

关键在于，新类型改进了设计模式的实现。例如，你现在可以用更简单的方式，对彼此无关的类型层级使用访问者（visitor）、模式匹配以及运行时多态。

本章你将学习：

\begin{itemize}
\item union 可能遇到哪些问题
\item 什么是带标签的联合体（discriminated union），以及为什么 union 需要类型安全
\item \texttt{std::variant} 如何工作、做了什么
\item 对 \texttt{std::variant} 的操作
\item 性能开销与内存需求
\item 使用示例
\end{itemize}

\section{基础}

在客户端代码中很少使用 union，大多数时候应当避免使用它。例如，在浮点运算中有一个“常见”的技巧：

\begin{codebox}[cpp]
union SuperFloat
{
float f;
int i;
}
int RawMantissa(SuperFloat f)
{
return f.i & ((1 << 23) - 1);
}
int RawExponent(SuperFloat f)
{
return (f.i >> 23) & 0xFF;
}
\end{codebox}

然而，虽然上面的代码在 C99 中可能有效，但由于更严格的别名规则（aliasing rules），它在 C++ 中是未定义行为！

关于这一点，现有的核心指南（Core Guidelines）有一条规则：

\texttt{C.183}：不要用 union 做类型双关（type punning）

原因：用与写入时不同的类型去读取 union 成员是未定义行为。这种双关是隐形的，至少比使用具名转型更难察觉。用 union 做类型双关是错误的一大来源。

在 C++17 之前，如果你想要类型安全的 union，可以使用 \texttt{boost::variant} 或其它第三方库。但现在你有了 \texttt{std::variant}。

下面是一个展示这个新类型能做什么的基础示例：

\begin{codebox}[cpp]
// Chapter Variant/variant_demo.cpp
#include <string>
#include <iostream>
#include <variant>
using namespace std;
// 用于打印当前激活的类型
struct PrintVisitor
{
void operator()(int i) { cout << "int: " << i << '\n'; }
void operator()(float f) { cout << "float: " << f << '\n'; }
void operator()(const string& s) { cout << "str: " << s << '\n'; }
};
int main()
{
variant<int, float, string> intFloatString;
static_assert(variant_size_v<decltype(intFloatString)> == 3);
// 默认初始化为第一个候选项，应为 0
visit(PrintVisitor{}, intFloatString);
// index 会显示当前使用的“类型”
cout << "index = " << intFloatString.index() << endl;
intFloatString = 100.0f;
cout << "index = " << intFloatString.index() << endl;
intFloatString = "hello super world";
cout << "index = " << intFloatString.index() << endl;
// 尝试使用 get_if：
if (const auto intPtr = get_if<int>(&intFloatString))
cout << "int: " << *intPtr << '\n';
else if (const auto floatPtr = get_if<float>(&intFloatString))
cout << "float: " << *floatPtr << '\n';
if (holds_alternative<int>(intFloatString))
cout << "the variant holds an int!\n";
else if (holds_alternative<float>(intFloatString))
cout << "the variant holds a float\n";
else if (holds_alternative<string>(intFloatString))
cout << "the variant holds a string\n";
// try/catch 与 bad_variant_access
try
{
auto f = get<float>(intFloatString);
cout << "float! " << f << '\n';
}
catch (bad_variant_access&)
{
cout << "our variant doesn't hold float at this moment...\n";
}
}
\end{codebox}

程序将输出：

\begin{codebox}
int: 0
index = 0
index = 1
index = 2
the variant holds a string
our variant doesn't hold float at this moment...
\end{codebox}

上面这个例子中有几个要点值得探讨：

\begin{itemize}
\item 第 18、22 行：如果你没有用一个值来初始化 variant，那么它会用第一个类型来初始化。这种情况下，第一个候选项类型必须拥有默认构造函数。第 22 行会打印出值 0。
\item 第 25、27、29、31、33、35 行：你可以通过 \texttt{index()} 或 \texttt{holds\_alternative} 检查当前使用的类型。
\item 第 39、41、47 行：你可以通过 \texttt{get\_if} 或 \texttt{get} 来访问值（可能会抛出 \texttt{bad\_variant\_access} 异常）。
\item 类型安全——variant 不允许你获取当前未激活类型的值。
\item 不会发生额外的堆分配。
\item 第 9、22 行：你可以使用访问者来调用当前激活类型的某个操作。示例中使用了 \texttt{PrintVisitor} 来打印当前激活的值。它是一个带有 \texttt{operator()} 重载的简单结构体。然后这个访问者被传给 \texttt{std::visit} 来执行访问。
\item variant 类会调用非平凡类型的析构函数和构造函数，因此在示例中，在我们切换到新的变体之前，\texttt{string} 对象会被清理掉。
\end{itemize}

\subsection{何时使用}

除非你正在做一些底层的、可能只涉及简单类型的工作，否则 union 也许是个可行的选择\footnote{参见 C++ Core Guidelines - Unions，了解 union 有效用法的示例。}。但对于所有其它需要可替代类型的用例，\texttt{std::variant} 才是正道。

一些可能的用途：

\begin{itemize}
\item 所有那些单个字段可能出现几种类型的地方：比如解析命令行、ini 文件、语言解析器等。
\item 高效地表达一次计算的若干种可能结果：例如求方程的根。
\item 错误处理——例如你可以返回 \texttt{variant<Object, ErrorCode>}。如果值可用，就返回 Object，否则赋一个错误码。
\item 有限状态机。
\item 无需虚表（vtable）和继承的多态（借助访问者模式）。
\end{itemize}

\subsection{函数式背景}

同样值得一提的是，variant 类型（也称为带标签的联合体、可辨识联合体或和类型，sum type）源自函数式语言世界与类型论\footnote{\url{https://en.wikipedia.org/wiki/Algebraic_data_type}}。

\section{创建 std::variant}

有几种方式可以创建并初始化 \texttt{std::variant}：

\begin{codebox}[cpp]
// Chapter Variant/variant_creation.cpp
// 默认初始化：（第一个类型必须有默认构造函数）
std::variant<int, float> intFloat;
std::cout << intFloat.index() << ", val: " << std::get<int>(intFloat) << '\n';
// 用 monostate 做默认初始化：
class NotSimple
{
public:
NotSimple(int, float) { }
};
// std::variant<NotSimple, int> cannotInit; // error
std::variant<std::monostate, NotSimple, int> okInit;
std::cout << okInit.index() << '\n';
// 传入一个值：
std::variant<int, float, std::string> intFloatString { 10.5f };
std::cout << intFloatString.index()
<< ", value " << std::get<float>(intFloatString) << '\n';
// 歧义
// double 可能转换为 float 或 int，因此编译器无法决定
//std::variant<int, float, std::string> intFloatString { 10.5 };
// 用 in_place 解决歧义
variant<long, float, std::string> longFloatString {
std::in_place_index<1>, 7.6 // double!
};
std::cout << longFloatString.index() << ", value "
<< std::get<float>(longFloatString) << '\n';
// 为复杂类型使用 in_place
std::variant<std::vector<int>, std::string> vecStr {
std::in_place_index<0>, { 0, 1, 2, 3 }
};
std::cout << vecStr.index() << ", vector size "
<< std::get<std::vector<int>>(vecStr).size() << '\n';
// 从另一个 variant 拷贝初始化：
std::variant<int, float> intFloatSecond { intFloat };
std::cout << intFloatSecond.index() << ", value "
<< std::get<int>(intFloatSecond) << '\n';
\end{codebox}

在这里把代码玩一玩 @Coliru\footnote{\url{http://coliru.stacked-crooked.com/a/1b930be2da3bd969}}。

\begin{itemize}
\item 默认情况下，variant 对象会用第一个类型初始化，
  \begin{itemize}
  \item 如果该类型没有默认构造函数，从而无法做到，那么你会得到一个编译错误
  \item 这种情况下你可以使用 \texttt{std::monostate} 作为第一个类型传进去
  \item monostate 让你能构建出“无值”的 variant，从而表现得类似 \texttt{optional}
  \end{itemize}
\item 你可以用一个值来初始化它，此时会使用最匹配的类型
  \begin{itemize}
  \item 如果存在歧义，那么你可以使用 \texttt{std::in\_place\_index} 版本来显式指明应当使用哪种类型
  \end{itemize}
\item \texttt{std::in\_place} 还允许你创建更复杂的类型，并向构造函数传入更多参数
\end{itemize}

\subsection{关于 std::monostate}

在示例中，你可能注意到了一个名为 \texttt{std::monostate} 的特殊类型。这是一个空类型，可以配合 variant 用来表示空状态。当第一个候选项没有默认构造函数时，这个类型会很方便。在这种情况下，你可以把 \texttt{std::monostate} 放在第一位（或者你也可以重排类型，找出一个有默认构造函数的类型）。

\subsection{就地构造}

\texttt{std::variant} 有两个可以使用的 in\_place 辅助工具：

\begin{itemize}
\item \texttt{std::in\_place\_type}——用于指定想要在 variant 中改变/设置的类型
\item \texttt{std::in\_place\_index}——用于指定想要改变/设置的索引。类型从 0 开始编号。
  \begin{itemize}
  \item 在 variant \texttt{std::variant<int, float, std::string>} 中——\texttt{int} 的索引是 0，\texttt{float} 的索引是 1，\texttt{string} 的索引是 2。这个索引与 \texttt{variant::index} 方法返回的值相同。
  \end{itemize}
\end{itemize}

好在，你并不总需要使用这些辅助工具来创建 variant。它足够智能，能够识别是否可以从传入的单个参数构造：

\begin{codebox}[cpp]
// 这会构造第二个/float：
std::variant<int, float, std::string> intFloatString { 10.5f };
\end{codebox}

对于 variant，我们至少在两种情况下需要这些辅助工具：

\begin{itemize}
\item 歧义——用来区分在多个可匹配类型中应当创建哪一个
\item 高效地创建复杂类型（与 \texttt{optional} 类似）
\end{itemize}

\textbf{歧义}

如果你写了这样的初始化：

\begin{codebox}[cpp]
std::variant<int, float> intFloat { 10.5 }; // 从 double 转换？
\end{codebox}

值 10.5 既可以被转换为 \texttt{int}，也可以被转换为 \texttt{float}，编译器不知道该应用哪种转换。不过它可能会报好几页的编译错误。

但你可以通过指明想要创建的类型来轻松处理这类错误：

\begin{codebox}[cpp]
std::variant<int, float> intFloat { std::in_place_index<0>, 10.5f };
// 或
std::variant<int, float> intFloat { std::in_place_type<int>, 10.5f };
\end{codebox}

\textbf{复杂类型}

与 \texttt{std::optional} 类似，如果你想高效地创建需要多个构造参数的对象，那就使用 \texttt{std::in\_place*}：

例如：

\begin{codebox}[cpp]
std::variant<std::vector<int>, std::string> vecStr {
std::in_place_index<0>, { 0, 1, 2, 3 } // 初始化列表被传入 vector
};
\end{codebox}

\section{修改值}

有四种方式可以修改 variant 的当前值：

\begin{itemize}
\item 赋值运算符
\item \texttt{emplace}
\item \texttt{get} 然后为当前激活的类型赋一个新值
\item 访问者（你不能改变类型，但可以改变当前候选项的值）
\end{itemize}

重要的是要知道，一切都是类型安全的，而且对象的生命期也得到了尊重。

\begin{codebox}[cpp]
// Chapter Variant/variant_changing_values.cpp
std::variant<int, float, std::string> intFloatString { "Hello" };
intFloatString = 10; // 现在我们是一个 int
intFloatString.emplace<2>(std::string("Hello")); // 现在我们又是一个 string
// std::get 返回引用，因此你可以修改其值：
std::get<std::string>(intFloatString) += std::string(" World");
intFloatString = 10.1f;
if (auto pFloat = std::get_if<float>(&intFloatString); pFloat)
*pFloat *= 2.0f;
\end{codebox}

\subsection{对象生命期}

当你使用 union 时，你需要管理内部状态：调用构造函数或析构函数。这很容易出错，也很容易搬起石头砸自己的脚。但 \texttt{std::variant} 会按你的预期管理对象生命期。这意味着，当它准备改变当前存储的类型时，会调用底层类型的析构函数。

\begin{codebox}[cpp]
std::variant<std::string, int> v { "Hello A Quite Long String" };
// v 为 string 分配了一些内存
v = 10; // 我们调用了 string 的析构函数！
// 没有内存泄漏
\end{codebox}

或者看这个带有自定义类型的示例：

\begin{codebox}[cpp]
// Chapter Variant/variant_lifetime.cpp
class MyType
{
public:
MyType() { std::cout << "MyType::MyType\n"; }
~MyType() { std::cout << "MyType::~MyType\n"; }
};
class OtherType
{
public:
OtherType() { std::cout << "OtherType::OtherType\n"; }
~OtherType() { std::cout << "OtherType::~OtherType\n"; }
};
int main()
{
std::variant<MyType, OtherType> v;
v = OtherType();
return 0;
}
\end{codebox}

这将产生如下输出：

\begin{codebox}
MyType::MyType
OtherType::OtherType
MyType::~MyType
OtherType::~OtherType
OtherType::~OtherType
\end{codebox}

在最开始，我们用 \texttt{MyType} 类型的默认值初始化；然后我们用一个 \texttt{OtherType} 实例来改变这个值，在赋值之前，会调用 \texttt{MyType} 的析构函数。随后我们销毁临时对象以及存储在 variant 中的对象。

\section{访问存储的值}

从你到目前为止看到的所有示例来看，你可能已经对如何访问值有了概念。但让我们对这一重要操作做个总结。

首先，即使你知道当前激活的是哪种类型，你也不能这样做：

\begin{codebox}[cpp]
std::variant<int, float, std::string> intFloatString { "Hello" };
std::string s = intFloatString;
// error: conversion from
// 'std::variant<int, float, std::string>'
// to non-scalar type 'std::string' requested
// std::string s = intFloatString;
\end{codebox}

所以你必须使用辅助函数来访问值。

你可以使用 \texttt{std::get<Type|Index>(variant)}，它是一个非成员函数。如果目标类型处于激活状态，它会返回该类型的引用（你可以传入类型或索引）。否则你会得到 \texttt{std::bad\_variant\_access} 异常。

\begin{codebox}[cpp]
std::variant<int, float, std::string> intFloatString;
try
{
auto f = std::get<float>(intFloatString);
std::cout << "float! " << f << '\n';
}
catch (std::bad_variant_access&)
{
std::cout << "our variant doesn't hold float at this moment...\n";
}
\end{codebox}

下一个选项是 \texttt{std::get\_if}。这个函数也是非成员的，并且不会抛出异常。它返回指向激活类型的指针，或者 \texttt{nullptr}。\texttt{std::get} 需要 variant 的引用，而 \texttt{std::get\_if} 接受指针。

\begin{codebox}[cpp]
if (const auto intPtr = std::get_if<0>(&intFloatString))
std::cout << "int!" << *intPtr << '\n';
\end{codebox}

不过，访问 variant 内部值最重要的方式大概还是使用访问者。

\section{用于 std::variant 的访问者}

随着 \texttt{std::variant} 的引入，我们还得到了一个便捷的 STL 函数 \texttt{std::visit}。它可以在所有传入的 variant 上调用给定的“访问者”。

声明如下：

\begin{codebox}[cpp]
template <class Visitor, class... Variants>
constexpr visit(Visitor&& vis, Variants&&... vars);
\end{codebox}

它会在 variant 的当前激活类型上调用 \texttt{vis}。

如果你只传入一个 variant，那么你必须为该 variant 中的类型准备好相应的重载。如果你传入两个 variant，那么你必须为你所传入 variant 的所有可能类型组合准备好重载。

访问者就是“一个可调用对象，它接受来自每个 variant 的每一种可能候选项”。

来看一些示例：

\begin{codebox}[cpp]
// 一个泛型 lambda：
auto PrintVisitor = [](const auto& t) { std::cout << t << '\n'; };
std::variant<int, float, std::string> intFloatString { "Hello" };
std::visit(PrintVisitor, intFloatString);
\end{codebox}

在上面的示例中，使用了泛型 lambda 来生成所有可能的重载。由于 variant 中的所有类型都支持 \texttt{<<}（流输出运算符），我们就能把它们打印出来。

在另一个例子中，我们可以用访问者来改变值：

\begin{codebox}[cpp]
auto PrintVisitor = [](const auto& t) { std::cout << t << '\n'; };
auto TwiceMoreVisitor = [](auto& t) { t*= 2; };
std::variant<int, float> intFloat { 20.4f };
std::visit(PrintVisitor, intFloat);
std::visit(TwiceMoreVisitor, intFloat);
std::visit(PrintVisitor, intFloat);
\end{codebox}

泛型 lambda 只有在我们的类型共享相同的“接口”时才能工作，但在大多数情况下，我们希望根据激活的类型执行不同的操作。

这就是为什么我们可以定义一个带有多个 \texttt{operator()} 重载的结构体：

\begin{codebox}[cpp]
struct MultiplyVisitor
{
float mFactor;
MultiplyVisitor(float factor) : mFactor(factor) { }
void operator()(int& i) const {
i *= static_cast<int>(mFactor);
}
void operator()(float& f) const {
f *= mFactor;
}
void operator()(std::string& ) const {
// nothing to do here...
}
};
std::visit(MultiplyVisitor(0.5f), intFloat);
std::visit(PrintVisitor, intFloat);
\end{codebox}

在示例中你会注意到，\texttt{MultiplyVisitor} 使用了一个状态来保存期望的缩放因子值。这为访问提供了很多可能性。

使用 lambda 时，我们习惯了把东西直接声明在使用的位置附近。而当你需要编写一个独立的结构体时，就得离开那个局部作用域。这就是为什么使用 overload 构造会很方便。

\subsection{Overload}

借助这个工具，你可以把所有匹配类型的 lambda 写在一个地方：

\begin{codebox}[cpp]
std::variant<int, float, std::string> myVariant;
std::visit(
overload {
[](const int& i) { std::cout << "int: " << i; },
[](const std::string& s) { std::cout << "string: " << s; },
[](const float& f) { std::cout << "float: " << f; }
},
myVariant
);
\end{codebox}

目前这个辅助工具还不是标准库的一部分（可能会在 C++20 中加入）。你可以用下面的代码实现它：

\begin{codebox}[cpp]
template<class... Ts> struct overload : Ts... { using Ts::operator()...; };
template<class... Ts> overload(Ts...) -> overload<Ts...>;
\end{codebox}

这段代码创建了一个继承自各个 lambda 的结构体，并使用它们的 \texttt{Ts::operator()}。整个结构体现在可以被传给 \texttt{std::visit}——它会从中选出合适的重载。

overload 用到了三个 C++17 特性：

\begin{itemize}
\item using 声明中的包展开（pack expansions）——配合可变参数模板的简短紧凑语法。
\item 自定义的模板实参推导规则——这让编译器能够推导出作为该模式基类的 lambda 的类型。没有它的话，我们就得定义一个“make”函数。
\item 对聚合初始化（aggregate initialisation）的扩展——overload 模式使用聚合初始化来初始化基类。在 C++17 之前这是不可能的。
\end{itemize}

这里是另一个使用它的示例：

\begin{codebox}[cpp]
// Chapter Variant/variant_overload.cpp
std::variant<int, float, std::string> intFloatString { "Hello" };
std::visit(overload{
[](int& i) { i*= 2; },
[](float& f) { f*= 2.0f; },
[](std::string& s) { s = s + s; }
}, intFloatString);
std::visit(PrintVisitor, intFloatString);
// prints: "HelloHello"
\end{codebox}

关于 \texttt{std::overload} 提案的论文是：P0051 - C++ generic overload function\footnote{\url{https://wg21.link/p0051r3}}。

你还可以在 bfilipek.com 上的这篇博客文章中读到更多关于 overload 模式机制的细节：2 Lines Of Code and 3 C++17 Features - The overload Pattern\footnote{\url{https://www.bfilipek.com/2019/02/2lines3featuresoverload.html}}。

\subsection{访问多个 variant}

\texttt{std::visit} 允许你不仅访问一个 variant，还可以在同一次调用中访问多个。不过重要的是要知道，对于多个 variant，你必须实现接受与输入 variant 数量相同参数个数的函数重载。而且你必须提供所有可能的类型组合。

例如对于：

\begin{codebox}[cpp]
std::variant<int, float, char> v1 { 's' };
std::variant<int, float, char> v2 { 10 };
\end{codebox}

如果你对这两个 variant 调用 \texttt{std::visit}，就必须提供 9 个函数重载：

\begin{codebox}[cpp]
std::visit(overload{
[](int a, int b) { },
[](int a, float b) { },
[](int a, char b) { },
[](float a, int b) { },
[](float a, float b) { },
[](float a, char b) { },
[](char a, int b) { },
[](char a, float b) { },
[](char a, char b) { }
}, v1, v2);
\end{codebox}

如果你漏掉一个重载，编译器就会报错。

看看下面的示例，其中每个 variant 代表一种食材，我们想把其中两种组合在一起：

\begin{codebox}[cpp]
// Chapter Variant/visit_multiple.cpp
#include <iostream>
#include <variant>
template<class... Ts> struct overload : Ts... { using Ts::operator()...; };
template<class... Ts> overload(Ts...) -> overload<Ts...>;
struct Pizza { };
struct Chocolate { };
struct Salami { };
struct IceCream { };
int main() {
std::variant<Pizza, Chocolate, Salami, IceCream> firstIngredient { IceCream() };
std::variant<Pizza, Chocolate, Salami, IceCream> secondIngredient { Chocolate() };
std::visit(overload{
[](const Pizza& p, const Salami& s) {
std::cout << "here you have, Pizza with Salami!\n";
},
[](const Salami& s, const Pizza& p) {
std::cout << "here you have, Pizza with Salami!\n";
},
[](const Chocolate& c, const IceCream& i) {
std::cout << "Chocolate with IceCream!\n";
},
[](const IceCream& i, const Chocolate& c) {
std::cout << "IceCream with a bit of Chocolate!\n";
},
[](const auto& a, const auto& b) {
std::cout << "invalid composition...\n";
},
}, firstIngredient, secondIngredient);
return 0;
}
\end{codebox}

程序将输出：\texttt{IceCream with a bit of Chocolate!}

上面的代码使用了 overload，并且用的是多个 lambda，而不是一个带有 \texttt{operator()} 重载的独立结构体。

有趣的是，这个示例只实现了“有效”食材组合，而“其余情况”由泛型 lambda（来自 C++14）处理。

一个泛型 lambda \texttt{[](const auto\& a, const auto\& b) \{ \}} 等价于下面这个可调用类型：

\begin{codebox}[cpp]
class UnnamedUniqueClass // << 编译器特定的名字...
{
public:
template<typename T, typename U>
auto operator () (const T& a, const T& b) const {
}
};
\end{codebox}

示例中使用的泛型 lambda 会为所有剩余的食材类型提供函数重载，由于它是一个模板，在决定最佳可行函数时，它总是排在具象重载（带有具体类型的 lambda）之后。

\section{其它 std::variant 操作}

为了完整起见：

\begin{itemize}
\item 你可以比较同类型的两个 variant：
  \begin{itemize}
  \item 如果它们包含相同的激活候选项，则调用相应的比较运算符。
  \item 如果一个 variant 拥有“更早”的候选项，那么它就“小于”拥有下一个激活候选项的 variant。
  \end{itemize}
\item variant 是一种值类型，所以你可以移动它。
\item 也可以对 variant 使用 \texttt{std::hash}。
\end{itemize}

\section{异常安全保证}

到目前为止一切看起来都很顺利……但当在 variant 中创建候选项时发生异常会怎样？

例如：

\begin{codebox}[cpp]
// Chapter Variant/variant_valueless.cpp
class ThrowingClass
{
public:
explicit ThrowingClass(int i) { if (i == 0) throw int (10); }
operator int () { throw int(10); }
};
int main(int argc, char** argv)
{
std::variant<int, ThrowingClass> v;
// 修改值：
try
{
v = ThrowingClass(0);
}
catch (...)
{
std::cout << "catch(...)\n";
// 我们保留了旧状态！
std::cout << v.valueless_by_exception() << '\n';
std::cout << std::get<int>(v) << '\n';
}
// 在 emplace 内部
try
{
v.emplace<0>(ThrowingClass(10)); // 调用 operator int
}
catch (...)
{
std::cout << "catch(...)\n";
// 旧状态已被销毁，因此我们不处于无效状态！
std::cout << v.valueless_by_exception() << '\n';
}
return 0;
}
\end{codebox}

在第一种情况——使用赋值运算符——异常是在该类型的构造函数中抛出的。这发生在 variant 中旧值被替换之前，因此 variant 的状态没有改变。如你所见，我们仍然可以访问 \texttt{int} 并打印它。

然而在第二种情况——\texttt{emplace}——异常是在 variant 的旧状态被销毁之后抛出的。\texttt{emplace} 调用 \texttt{operator int} 来替换值，但它抛出了异常。此后，variant 处于错误的状态，我们无法恢复之前的状态。

还要注意，一个“因异常而无值（valueless by exception）”的 variant 处于无效状态。无法从这样的 variant 访问值。这就是为什么 \texttt{variant::index} 会返回 \texttt{variant\_npos}，而 \texttt{std::get} 和 \texttt{std::visit} 会抛出 \texttt{bad\_variant\_access}。

\section{性能与内存考量}

\texttt{std::variant} 使用和 union 类似的方式来占用内存：它会取底层类型的最大尺寸。但由于我们需要某种东西来知道当前激活的是哪个候选项，就需要多占用一些空间。此外，一切都需要遵循对齐规则。

下面是一些基本的大小：

\begin{codebox}[cpp]
// Chapter Variant/variant_sizeof.cpp
std::cout << "sizeof string: "
<< sizeof(std::string) << '\n';
std::cout << "sizeof variant<int, string>: "
<< sizeof(std::variant<int, std::string>) << '\n';
std::cout << "sizeof variant<int, float>: "
<< sizeof(std::variant<int, float>) << '\n';
std::cout << "sizeof variant<int, double>: "
<< sizeof(std::variant<int, double>) << '\n';
\end{codebox}

在 GCC 8.1、32 位下：

\begin{codebox}
sizeof string: 32
sizeof variant<int, string>: 40
sizeof variant<int, float>: 8
sizeof variant<int, double>: 16
\end{codebox}

更有趣的是，\texttt{std::variant} 不会分配任何额外的空间！不会为了持有 variant 或判别器而发生动态分配。

为了拥有一个安全的求和类型，你要付出内存占用增大的代价。多出来的比特可能会影响 CPU 缓存。因此，对你应用中使用 variant 的热点代码，你可能想做一些基准测试。

\section{从 boost::variant 迁移}

Boost Variant 大约在 2004 年引入，所以在 \texttt{std::variant} 加入标准之前，它已经有了 13 年的经验积累。STL 中的这个类型汲取了 boost 版本的经验并做了改进。

主要变化如下：

\begin{table}[H]
\centering
\begin{tabularx}{\textwidth}{l X X}
\toprule
特性 & Boost.Variant (1.67.0) & std::variant \\
\midrule
额外内存分配 & 赋值时可能发生（参见设计概述） & 从不 \\
空状态⁷ & 否 & 是 \\
访问（visiting） & \texttt{apply\_visitor} & \texttt{std::visit} \\
按下标访问 & 否 & 是 \\
递归 variant & 是（参见 \texttt{make\_recursive\_variant}⁸） & 否 \\
重复项 & 否 & 是 \\
空的可选项 & \texttt{boost::blank} & \texttt{std::monostate} \\
\bottomrule
\end{tabularx}
\end{table}

\noindent ⁷ \url{https://www.boost.org/doc/libs/1_67_0/doc/html/variant/design.html}

\noindent ⁸ \url{https://www.boost.org/doc/libs/1_67_0/doc/html/boost/make_recursive_variant.html}

\section{std::variant 示例}

在了解了 \texttt{std::variant} 的大部分细节之后，我们现在可以探究几个示例了。

\subsection{错误处理}

基本思想是：用某个 ErrorCode 把可能的返回类型包起来，从而让函数输出更多关于错误的信息。无需使用异常或输出参数。这类似于 \texttt{std::expected}——一个计划在未来 C++ 标准中引入的新类型。

\begin{codebox}[cpp]
// Chapter Variant/variant_error_handling.cpp
enum class ErrorCode
{
Ok,
SystemError,
IoError,
NetworkError
};
std::variant<std::string, ErrorCode> FetchNameFromNetwork(int i)
{
if (i == 0)
return ErrorCode::SystemError;
if (i == 1)
return ErrorCode::NetworkError;
return std::string("Hello World!");
}
int main()
{
auto response = FetchNameFromNetwork(0);
if (std::holds_alternative<std::string>(response))
std::cout << std::get<std::string>(response) << "\n";
else
std::cout << "Error!\n";
response = FetchNameFromNetwork(10);
if (std::holds_alternative<std::string>(response))
std::cout << std::get<std::string>(response) << "\n";
else
std::cout << "Error!\n";
return 0;
}
\end{codebox}

在示例中，返回的是 ErrorCode 或一个常规类型。

\subsection{解析命令行}

命令行可能包含可以被多种方式解释的文本参数：

\begin{itemize}
\item 作为整数
\item 作为浮点数
\item 作为布尔标志
\item 作为字符串（不解析）
\item ……
\end{itemize}

我们可以构建一个能持有所有可能选项的 variant。

下面是一个 \texttt{int}、\texttt{float} 和 \texttt{string} 的简单版本：

\begin{codebox}[cpp]
// Chapter Variant/variant_parsing_int_float.cpp
class CmdLine
{
public:
using Arg = std::variant<int, float, std::string>;
private:
std::map<std::string, Arg> mParsedArgs;
public:
explicit CmdLine(int argc, const char** argv) { ParseArgs(argc, argv); }
std::optional<Arg> Find(const std::string& name) const;
// ...
};
\end{codebox}

解析代码如下：

\begin{codebox}[cpp]
// Chapter Variant/variant_parsing_int_float.cpp
CmdLine::Arg TryParseString(std::string_view sv)
{
// 先尝试按 float 解析
float fResult = 0.0f;
const auto last = sv.data() + sv.size();
const auto res = std::from_chars(sv.data(), last, fResult);
if (res.ec != std::errc{} || res.ptr != last)
{
// 如果做不到，就假定它是 string
return std::string{sv};
}
// 没有小数部分？那就转换为整数
if (static_cast<int>(fResult) == fResult)
return static_cast<int>(fResult);
return fResult;
}
void CmdLine::ParseArgs(int argc, const char** argv)
{
// 形式：-argName value -argName value
for (int i = 1; i < argc; i+=2)
{
if (argv[i][0] != '-') // 超级高级的模式匹配！ :)
throw std::runtime_error("wrong command name");
mParsedArgs[argv[i]+1] = TryParseString(argv[i+1]);
}
}
\end{codebox}

在撰写本文时，GCC/Clang 中的 \texttt{std::from\_chars} 只支持整数。MSVC 2017 15.8 则对浮点数也有完整支持。你可以在独立的“字符串转换”一章中阅读更多关于 \texttt{from\_chars} 的内容。

\texttt{TryParseString} 的思路是尝试把输入字符串解析成最匹配的类型。所以如果它看起来像整数，我们就尝试取一个整数。否则，我们就返回一个未解析的字符串。当然，我们可以扩展这种做法。

解析完成后，客户端可以使用 \texttt{Find()} 方法来检查某个参数是否存在：

\begin{codebox}[cpp]
std::optional<CmdLine::Arg> CmdLine::Find(const std::string& name) const
{
if (const auto it = mParsedArgs.find(name); it != mParsedArgs.end())
return it->second;
return { };
}
\end{codebox}

\texttt{Find()} 使用 \texttt{std::optional} 来返回值。如果在 map 中找不到该参数，客户端就会得到一个空的 optional。

我们如何使用它的示例：

\begin{codebox}[cpp]
// Chapter Variant/variant_parsing_int_float.cpp
try
{
CmdLine cmdLine(argc, argv);
if (auto arg = cmdLine.Find("paramInt"); arg)
std::cout << "paramInt is " << std::get<int>(*arg) << '\n';
if (auto arg = cmdLine.Find("paramFloat"); arg)
{
if (const auto intPtr = std::get_if<int>(&*arg); intPtr)
std::cout << "paramFloat is " << *intPtr << " (integer)\n";
else
std::cout << "paramFloat is " << std::get<float>(*arg) << '\n';
}
if (auto arg = cmdLine.Find("paramText"); arg)
std::cout << "paramText is " << std::get<std::string>(*arg) << '\n';
}
catch (std::bad_variant_access& err)
{
std::cerr << " ...err: accessing a wrong variant type, "
<< err.what() << '\n';
}
catch (std::runtime_error &err)
{
std::cerr << " ...err: " << err.what() << '\n';
}
\end{codebox}

上面的示例使用 \texttt{cmdLine.Find()} 来检查是否存在给定参数。它返回 \texttt{std::optional}，因此我们必须检查它是否为空。当我们确定参数可用时，就可以检查它的类型。

\texttt{CmdLine} 会尝试找到最匹配的类型，例如对于浮点数，我们存在歧义——因为 90 也是 float，但代码会把它存为 int（因为它没有小数部分）。

要解决这类歧义，我们可以传入一些关于期望类型的额外信息，或者提供某些辅助方法。

\subsection{解析配置文件}

这个想法来自前面命令行的示例。在配置文件的情况下，我们通常处理的是 \texttt{<Name, Value>} 这样的键值对。其中 Value 可能是不同的类型：\texttt{string}、\texttt{int}、\texttt{array}、\texttt{bool}、\texttt{float} 等等。

对于这样的用例，甚至可以用 \texttt{void*} 来持有这种未知类型。但这种模式极其容易出错。如果我们已知所有可能的类型，就可以用 \texttt{std::variant} 来改进设计，或者借助 \texttt{std::any}。

\subsection{状态机}

如何建模一个状态机？例如门的状态：

\myGraphic{0.9}{images/fig\_p133\_1.png}{（门状态机示意图）}

我们可以使用不同类型的状态，并把访问者用作事件：

\begin{codebox}[cpp]
// Chapter Variant/variant_fsm.cpp
struct DoorState
{
struct DoorOpened {};
struct DoorClosed {};
struct DoorLocked {};
using State = std::variant<DoorOpened, DoorClosed, DoorLocked>;
void open()
{
m_state = std::visit(OpenEvent{}, m_state);
}
void close()
{
m_state = std::visit(CloseEvent{}, m_state);
}
void lock()
{
m_state = std::visit(LockEvent{}, m_state);
}
void unlock()
{
m_state = std::visit(UnlockEvent{}, m_state);
}
State m_state;
};
\end{codebox}

这里是相应的事件：

\begin{codebox}[cpp]
// Chapter Variant/variant_fsm.cpp
struct OpenEvent
{
State operator()(const DoorOpened&){ return DoorOpened(); }
State operator()(const DoorClosed&){ return DoorOpened(); }
// 无法打开已锁上的门
State operator()(const DoorLocked&){ return DoorLocked(); }
};
struct CloseEvent
{
State operator()(const DoorOpened&){ return DoorClosed(); }
State operator()(const DoorClosed&){ return DoorClosed(); }
State operator()(const DoorLocked&){ return DoorLocked(); }
};
struct LockEvent
{
// 无法锁上已打开的门
State operator()(const DoorOpened&){ return DoorOpened(); }
State operator()(const DoorClosed&){ return DoorLocked(); }
State operator()(const DoorLocked&){ return DoorLocked(); }
};
struct UnlockEvent
{
// 无法解锁已打开的门
State operator()(const DoorOpened&){ return DoorOpened(); }
State operator()(const DoorClosed&){ return DoorClosed(); }
// 解锁
State operator()(const DoorLocked&){ return DoorClosed(); }
};
\end{codebox}

\subsection{多态}

在 C++ 中，大多数时候我们可以安全地使用基于虚表（vtable）方式的运行时多态。你有一组共享相同接口的相关类型，并且有一个定义良好的虚函数可供调用。

但如果你有“互不相干”的类型，它们并不共享同一个基类，那该怎么办？如果你希望快速添加新功能，又不想改动所支持类型的代码，那又该怎么办？

借助 \texttt{std::variant} 和 \texttt{std::visit}，我们可以构建如下示例：

\begin{codebox}[cpp]
// Chapter Variant/variant_polymorphism.cpp
class Triangle
{
public:
void Render()
{ std::cout << "Drawing a triangle!\n"; }
};
class Polygon
{
public:
void Render() { std::cout << "Drawing a polygon!\n"; }
};
class Sphere
{
public:
void Render() { std::cout << "Drawing a sphere!\n"; }
};
int main()
{
std::vector<std::variant<Triangle, Polygon, Sphere>> objects {
Polygon(),
Triangle(),
Sphere(),
Triangle()
};
auto CallRender = [](auto& obj) { obj.Render(); };
for (auto& obj : objects)
std::visit(CallRender, obj);
}
\end{codebox}

上面的示例只展示了从互不相干的类型调用方法的第一种情况。它把所有可能出现的形状类型包装进一个 variant，然后使用访问者把调用分派到正确的类型。

例如，如果你想对对象排序，那么可以编写另一个持有某些状态的访问者。这样一来，无需改动这些类型，你就能获得更多功能。
